\documentclass[11pt]{article}

\usepackage[margin=1in]{geometry}
\usepackage{amsmath, amssymb, amsthm}
\usepackage{mathtools}
\usepackage[colorlinks=true, linkcolor=blue, urlcolor=blue, citecolor=blue]{hyperref}

\title{\textbf{Keep AI calls out of Convex mutations}}
\author{Sanjay Suresh}
\date{March 2026}

\begin{document}
\maketitle

Convex mutations are transactional and meant to be deterministic --- they can re-run, and the
runtime assumes the same inputs produce the same writes. Anything non-deterministic or network-bound
(an LLM call, token streaming, an embeddings request) breaks that assumption.

The right shape is: an \texttt{action} does the external call, then hands its result to a
\texttt{mutation} that persists it. The action is where side effects live; the mutation stays a
clean, replayable write.

I learned this the slow way on AIManager --- chat responses were occasionally double-writing because
the model call sat inside the mutation and got retried. Moving the call into an action made the whole
thing predictable.

\end{document}
